\section{Tolerances}
\label{sec:theory_general_tolerances}

In real-world resistors, capacitors, and inductors (and actually
a lot of other components) cannot be manufactured to their exact
specified values. This is due to natural imperfections of
process, material, and environment.

To account for this, manufacturers tend to design for a approximate
value and specify a tolerance range. For example, a resistor
labeled as \(100 \Omega\) with a tolerance of \(\pm 5\%\)
means that the actual resistance value can vary between
\(95 \Omega\) and \(105 \Omega\) even though it is labeled as
\(100 \Omega\).


\subsection{Mathematical Tricks}
A really interesting mathematical trick you can do is measure a
bunch of 100 \(\Omega\) resistors that all say have a specifed
\(\pm 5\%\) tolerance. You'll find that the recorded data points
will follow a normal distribution (bell curve) centered around
\(100 \Omega\).

Manufacturers use statistical properties like Standard Deviation,
Cumulative Distribution Function (CDF), and Probability Density
Function (PDF) to assess the quality of their manufacturing
process and ensure that most of their components fall within
the specified tolerance range. This extends not just to resistors
but also capacitors, inductors, integrated circuits, and many
other electronic components.
